\begin{thema}{Δ}
Στο σχήμα δίνεται η $C_f$ μιας συνεχούς στο $[0,6]$ και δύο φορές παραγωγίσιμης στο $(0,6)$ συνάρτησης $f:[0, 6]\to\mathbb{R}$ (μπλε) μαζί με την ευθεία $y = x$ (διακεκομμένη). Η $C_f$ τέμνει τη $y = x$ ακριβώς στα $x = 0$, $x = 3$ και $x = 5$. Δίνεται ότι $f(2) = 3$, $f(4) = 3$ και $f(6)=\dfrac{13}{2}$.

\begin{center}
\begin{tikzpicture}[scale=0.7]
  \draw[help lines, gray!30] (-1,-1) grid (7,7);
  \draw[-latex] (-1.3,0) -- (7.3,0) node[right] {$x$};
  \draw[-latex] (0,-1.3) -- (0,7.3) node[above] {$y$};
  \foreach \x in {1,2,3,4,5,6}
    \draw (\x,0.1) -- (\x,-0.1) node[below,font=\footnotesize] {$\x$};
  \foreach \y in {1,2,3,4,5,6}
    \draw (0.1,\y) -- (-0.1,\y) node[left,font=\footnotesize] {$\y$};
  \draw[dashed, gray!50] (0,0) -- (6,6);
  \draw[very thick, maincolor, smooth, tension=0.5]
    plot coordinates {(0,0) (1,2) (2,3) (3,3) (4,3) (5,5) (6,6.5)};
  \filldraw[maincolor] (0,0) circle (2pt);
  \filldraw[maincolor] (2,3) circle (2pt) node[above left,font=\footnotesize] {$(2,3)$};
  \filldraw[maincolor] (3,3) circle (2pt) node[below right,font=\footnotesize] {$(3,3)$};
  \filldraw[maincolor] (5,5) circle (2pt) node[below right,font=\footnotesize] {$(5,5)$};
\end{tikzpicture}
\end{center}

\begin{erwthma}
  \item Να λύσετε την εξίσωση $f(x) = x$ και την ανίσωση $f(x) > x$ στο διάστημα $[0, 6]$.
  \item Να αποδείξετε ότι η εξίσωση $f'(x) = 1$ έχει τουλάχιστον δύο ρίζες στο διάστημα $(0, 5)$.
  \item Να αποδείξετε ότι υπάρχει $\eta \in (0, 5)$ τέτοιο ώστε $f''(\eta) = 0$.
  \item Να αποδείξετε ότι $\dint_0^3 f(x)\,dx > \frac{9}{2}$.
\end{erwthma}
\end{thema}
